%%%%%%%%%%%%%%%%%%%%%%%%%%%%%%%%%%%%%%%%%%%%%%%%%%%%%%%%%%%%%%%%%%%%%%%%%%%%%%%%
%2345678901234567890123456789012345678901234567890123456789012345678901234567890
%        1         2         3         4         5         6         7         8

\documentclass[letterpaper, 10 pt, conference]{ieeeconf}  % Comment this line out if you need a4paper

%\documentclass[a4paper, 10pt, conference]{ieeeconf}      % Use this line for a4 paper

\IEEEoverridecommandlockouts                              % This command is only needed if
                                                          % you want to use the \thanks command

\overrideIEEEmargins                                      % Needed to meet printer requirements.

%In case you encounter the following error:
%Error 1010 The PDF file may be corrupt (unable to open PDF file) OR
%Error 1000 An error occurred while parsing a contents stream. Unable to analyze the PDF file.
%This is a known problem with pdfLaTeX conversion filter. The file cannot be opened with acrobat reader
%Please use one of the alternatives below to circumvent this error by uncommenting one or the other
%\pdfobjcompresslevel=0
%\pdfminorversion=4

% The following packages can be found on http:\\www.ctan.org
%\usepackage{graphics} % for pdf, bitmapped graphics files
%\usepackage{epsfig} % for postscript graphics files
%\usepackage{mathptmx} % assumes new font selection scheme installed
%\usepackage{times} % assumes new font selection scheme installed
\usepackage{amsmath} % assumes amsmath package installed
\usepackage{amssymb}  % assumes amsmath package installed

\usepackage{graphicx}
\usepackage{booktabs}
\usepackage{balance}

% arXiv includes the supplement; set false for the eight-page RA-L version.
\newif\ifincludesupplementary
\includesupplementaryfalse

% Compact float layout; keep the class's text size and page margins.
\setlength{\textfloatsep}{10pt plus 2pt minus 2pt}
\setlength{\dbltextfloatsep}{10pt plus 2pt minus 2pt}
\setlength{\floatsep}{8pt plus 2pt minus 2pt}
\setlength{\dblfloatsep}{8pt plus 2pt minus 2pt}
\setlength{\intextsep}{8pt plus 2pt minus 2pt}
\renewcommand{\textfraction}{0.08}
\renewcommand{\bottomfraction}{0.7}
\setcounter{topnumber}{4}
\setcounter{totalnumber}{6}

\makeatletter
\g@addto@macro\normalsize{%
  \setlength{\abovedisplayskip}{6pt plus 2pt minus 2pt}%
  \setlength{\belowdisplayskip}{6pt plus 2pt minus 2pt}%
  \setlength{\abovedisplayshortskip}{0pt plus 1pt}%
  \setlength{\belowdisplayshortskip}{4pt plus 1pt minus 1pt}%
}
\makeatother

\title{\LARGE \bf
GEAR-Nav: Geometric Energy-Guided Aggregation and Refinement for Vision-Language Navigation
}


\author{anonymous
}


\makeatletter
\IEEEaftertitletext{%
  \begin{minipage}{\textwidth}
    \centering
    \setlength{\abovecaptionskip}{4pt}
    \includegraphics[width=\linewidth]{fig/astar_teaser.pdf}
    \vspace{-10pt}% Tighten the class's non-floating image-to-caption spacing.
    \def\@captype{figure}
    \caption[GEAR learned trajectory refinement]{
      \textbf{Learning to refine paths from geometric energy.}
      GEAR learns a trajectory-correction field from energy-derived supervision,
      then repeatedly refines a candidate path using the image, goal, and current path.
      Across (a) a caf\'{e} sidewalk, (b) a rainy night sidewalk, (c) an indoor dining
      concourse, and (d) a market passage, conceptual stages extend and reshape short
      sampled priors toward the goal. Later paths are authored to illustrate progression
      observed in a recorded rollout; stage numbers are not update counts. Each ESDF
      and uncropped camera view shows the same metric paths through the saved camera
      projection. ESDFs visualize training geometry; inference requires neither the
      distance field nor the energy function.
    }
    \label{fig:teaser}
  \end{minipage}\vskip0.5\baselineskip%
}
\makeatother


\begin{document}
\typeout{REVIEW_LAYOUT: textwidth=\the\textwidth; columnwidth=\the\columnwidth; textheight=\the\textheight}



\maketitle
\thispagestyle{empty}
\pagestyle{empty}


%%%%%%%%%%%%%%%%%%%%%%%%%%%%%%%%%%%%%%%%%%%%%%%%%%%%%%%%%%%%%%%%%%%%%%%%%%%%%%%%
\begin{abstract}

Vision-Language-Action (VLA) models provide a promising framework for goal-conditioned robot navigation, but training obstacle-aware policies typically relies on target trajectories obtained from robot demonstrations or offline trajectory optimization. We introduce GEAR-Nav (GEAR), a framework for training navigation VLAs from in-the-wild egocentric video without target trajectory supervision. Given an RGB observation and navigation goal, we reconstruct local scene geometry and define an analytic energy over candidate waypoint trajectories that captures goal reaching, collision avoidance, progress, and motion smoothness. Rather than solving this energy to produce a trajectory label, we use its local descent direction as privileged supervision for learning an autonomous trajectory-refinement field. Starting from simple, kinematically feasible trajectory priors, the learned field iteratively improves candidate paths. To address the distribution shift between initial priors and model-induced trajectories, we introduce an energy-oracle dataset aggregation procedure that repeatedly queries projection-aware energy-descent directions at trajectories visited by the current model, aggregates these state-direction pairs, and trains the VLA on the resulting replay distribution. Geometry and the analytic energy are used only during training; at inference time, the model iteratively refines trajectories using visual observations, goal conditioning, and kinematic projection alone, without access to a distance field, energy function, or online trajectory optimizer. We evaluate our approach on held-out egocentric scenes and real-world robot navigation tasks, measuring goal reaching, collision avoidance, obstacle clearance, and trajectory progress. Our results show that [RESULTS]. These findings demonstrate that geometric energy can directly supervise navigation VLAs, eliminating the need to first convert scene geometry into planner-generated target trajectories.

\end{abstract}


%%%%%%%%%%%%%%%%%%%%%%%%%%%%%%%%%%%%%%%%%%%%%%%%%%%%%%%%%%%%%%%%%%%%%%%%%%%%%%%%
\section{INTRODUCTION}
\label{sec:introduction}

Vision-Language-Action (VLA) models connect visual observations and
language instructions to robot actions, bringing semantic knowledge
from large-scale pretraining into embodied control~\cite{pi05}.
Navigation VLAs extend this interface to destinations specified by
language, images, or spatial coordinates~\cite{navila,omnivla}.
Scaling these capabilities to everyday mobile robots, however, requires
supervision that covers diverse environments, goals, and obstacle
configurations. Collecting robot demonstrations across these combinations
is costly, particularly in cluttered indoor and outdoor spaces where
recognizing a destination must be coupled with close-range geometric
reasoning.

The challenge is not simply to record more motion. A traversal shows how
to move toward one destination, while the same observation may support
many other goals, each requiring a different route through the available
free space. Reaching a doorway, approaching an object, or passing around
furniture can demand different choices even from the same starting
viewpoint. A broadly goal-conditioned policy therefore needs supervision
that relates the requested destination to obstacle clearance,
traversability, and progress, rather than relying only on the path that
happened to be demonstrated. Geometric planning objectives express these
requirements explicitly; transferring them into a visual policy offers a
way to teach obstacle-aware behavior without requiring online geometric
mapping at deployment.

In-the-wild egocentric videos offer a scalable source of observations for
this purpose. Walking tours and other first-person recordings capture
varied layouts, viewpoints, and obstacles without dedicated robot data
collection. Yet these videos do not directly provide robot-executable
paths toward every potential goal in a scene. Existing approaches recover
motion from video~\cite{citywalker} or use a pretrained navigation policy
to supply trajectory supervision~\cite{lelan}. Reconstructed geometry
provides another route: VEGA~\cite{vega} generates obstacle-aware reference
paths toward scene-derived goals and uses them to train a flow-matching
navigation VLA. In each case, learning relies on a target trajectory,
whether observed, predicted, or planned. Geometry, however, contains more
than a preferred final path: it also indicates how a candidate path
should change to improve clearance and goal-directed progress.

This raises the question: \emph{can geometric objectives teach a navigation
VLA to refine trajectories without first producing target paths?}
Learning navigation directly from differentiable costs is established by
methods such as iPlanner and ViPlanner~\cite{iplanner,viplanner}.
Our focus is on learning local corrections conditioned on the image,
goal, and current candidate trajectory, so that the same model can be
applied repeatedly as that trajectory evolves. This requires both a
geometric supervision signal that respects the constraints on path
updates and training coverage of the candidate paths encountered during
learned refinement.

We introduce \textbf{GEAR-Nav} (GEAR), \emph{Geometric Energy-Guided
Aggregation and Refinement for Vision-Language Navigation}. GEAR adapts a
pretrained VLA using egocentric observations and goals expressed through
language, image regions, or robot-frame waypoints. During training,
reconstructed scene geometry defines an analytic energy over candidate
paths that encodes collision avoidance, obstacle clearance, goal reaching,
progress, and smooth motion. Instead of minimizing this energy to create
a reference trajectory, we query its local descent direction. A
projection-aware oracle scales this correction through backtracking on
projected proposals, accounting for the fixed origin and bounded waypoint
segments. The VLA learns to predict the resulting correction from the
observation, goal, and candidate path. Its field is autonomous: the update
rule depends on the current path, without a refinement-time input or a
noise-to-target interpolation.

Repeated prediction creates a second challenge: each learned correction
changes the candidate supplied to the next update. Training only on
initial paths leaves these later inputs outside the supervised
distribution. Inspired by dataset aggregation~\cite{dagger}, GEAR starts
from goal-biased unicycle priors, advances candidate paths with the
current learned field, queries oracle corrections at the resulting
states, and retains these examples across training rounds. An energy
acceptance guard controls frontier updates during collection. At
deployment, both the guard and the oracle are removed: refinement uses
only the RGB observation, goal, candidate path, and lightweight kinematic
projection. Neither an ESDF nor analytic energy evaluations or gradients
are required at inference. The absence of target trajectories applies to
this navigation adaptation; it does not imply action-free pretraining of
the VLA backbone.

Our contributions are threefold:
\begin{itemize}
    \item \textbf{Local geometric supervision for navigation VLAs:}
    We formulate navigation adaptation using energy-derived trajectory
    corrections, without demonstrated or planner-generated target paths.

    \item \textbf{Projection-aware supervision and trajectory-space
    aggregation:} We construct local correction labels consistent with
    projected updates and aggregate them at model-induced candidate paths
    to train repeated refinement beyond the initial prior distribution.

    \item \textbf{Autonomous visual trajectory refinement:} We adapt a
    pretrained VLA into a reusable, multimodal goal-conditioned correction
    field that refines candidate paths without the training-time geometry,
    analytic energy, or acceptance guard at deployment.
\end{itemize}

% Add a concise, evidence-backed experimental summary here once
% results are finalized.

%%%%%%%%%%%%%%%%%%%%%%%%%%%%%%%%%%%%%%%%%%%%%%%%%%%%%%%%%%%%%%%%%%%%%%%%%%%%%%%%
\section{RELATED WORK}
\label{sec:related_work}

\subsection{Navigation VLAs and Learning from Video}

Navigation VLAs connect visual-language representations with embodied
control. NaVILA~\cite{navila} uses a hierarchical architecture in which a
VLA generates intermediate spatial commands for a visual locomotion
policy. OmniVLA~\cite{omnivla} supports language, image, and spatial goals,
including combinations of modalities. These approaches establish useful
interfaces for goal-conditioned navigation; GEAR focuses on the
supervision used to learn obstacle-aware waypoint updates within such an
interface.

Egocentric video provides an alternative to collecting task-specific robot
demonstrations. CityWalker~\cite{citywalker} extracts motion supervision
from web-scale walking and driving videos using visual odometry.
LeLaN~\cite{lelan} combines semantic goal annotation and goal-reaching
objectives with trajectory-based distillation from a pretrained navigation
model to teach obstacle avoidance. VEGA~\cite{vega} instead reconstructs
local geometry and generates MPPI trajectories toward multiple multimodal
goals, using these paths to supervise a flow-matching VLA. GEAR shares
VEGA's use of video-derived geometry and aligned goals, but changes the
supervision: it queries local corrections at candidate paths instead of
constructing optimized target trajectories. Thus, its distinction is not
the use of action-free video or multimodal goals alone, but how geometry
supervises the navigation model.

\subsection{Cost-Supervised Planning and Learned Optimization}

Trajectory optimization methods such as CHOMP~\cite{chomp} improve initial
paths using geometric gradients and smoothness objectives. Their update
rules provide a classical foundation for iterative trajectory refinement,
but require access to the planning objective when solving a new problem.
Learning-based approaches can transfer these objectives into a policy.
iPlanner~\cite{iplanner} trains a navigation network through differentiable
planning costs without labeled trajectories, while ViPlanner~\cite{viplanner}
extends imperative learning to geometric and semantic traversability.
Consequently, learning navigation from costs without target paths is
established prior work.

GEAR differs in the quantity it learns: a conditional correction field
over candidate trajectories. Analytic gradients are converted into cached,
projection-aware direction labels, and network training regresses those
labels rather than backpropagating the trajectory energy through the
network. This also relates to learned optimization, where parameterized
update rules exploit structure shared across optimization
problems~\cite{learnedoptimizer}. Unlike the gradient-input recurrent
optimizer of Andrychowicz et al., GEAR predicts updates from visual
observations, goal conditioning, and the current path; analytic energy
gradients are training targets, not deployment inputs.

\subsection{Generative Trajectory Models and Cost Guidance}

Diffusion Policy~\cite{diffusionpolicy} models robot actions through
conditional denoising, while flow matching~\cite{flowmatching} learns
transport fields between source and target distributions. These methods
motivate expressive action generators, but their conventional supervised
formulations learn from action or trajectory samples. Motion Planning
Diffusion~\cite{mpd} combines a learned trajectory prior with cost
gradients during denoising, allowing generated motions to adapt to new
planning constraints. Such cost-guided inference retains explicit access
to the task objective.

GEAR also performs iterative updates, but its training targets are local
energy corrections rather than denoising targets or noise-to-trajectory
displacements. It learns an autonomous refinement field, with no
refinement-time conditioning, and removes the analytic energy at
deployment. Its objective is to learn repeated trajectory improvement,
not to match a prescribed trajectory distribution. This separates GEAR
from both trajectory-supervised generative policies and inference-time
cost-guided samplers.

\subsection{Dataset Aggregation and Planner Distillation}

DAgger~\cite{dagger} addresses sequential distribution shift by querying
expert labels at states encountered by the learner and aggregating these
examples across training iterations. This principle has also been applied
to planning: Deep-PANTHER~\cite{deeppanther} uses dataset aggregation to
train a trajectory-generating policy from an optimization-based expert.
These precedents distinguish dataset aggregation itself from the specific
oracle and state space used in GEAR.

In GEAR, the aggregated states are candidate waypoint paths associated
with fixed observation-goal conditions, and the labels are local
energy-descent directions rather than complete expert trajectories. The
learned field advances the candidate frontier under a training-only
energy guard, and subsequent oracle queries supervise corrections at the
new states. The resulting method combines trajectory-space aggregation
with projection-aware local supervision, while evaluating the deployed
policy without the oracle or acceptance guard.


%%%%%%%%%%%%%%%%%%%%%%%%%%%%%%%%%%%%%%%%%%%%%%%%%%%%%%%%%%%%%%%%%%%%%%%%%%%%%%%%
\section{BACKGROUND}
\label{sec:background}

\subsection{Goal-Conditioned Navigation VLAs}
\label{sec:bg_vla}

Let $I$ denote the current RGB observation and $g$ the goal
specification, which may be a language instruction, an image region,
or a robot-frame coordinate. We represent a local navigation
trajectory as
\begin{equation}
    x = (p_0,\ldots,p_{H-1}) \in \mathbb{R}^{H\times 2},
    \qquad p_0 = (0,0),
\end{equation}
where $H$ is the number of waypoints and all positions are expressed
in the robot's local frame. A navigation policy models
$\pi_\theta(x\mid I,g)$. Pretrained VLAs provide visual-language
context that can condition an action expert, while
navigation-specific encoders accommodate spatial goal
inputs~\cite{pi05,omnivla}. In this work, \emph{action-free}
refers to the task-specific training data: no demonstrated actions
or optimized target paths supervise GEAR's navigation adaptation.
It does not imply that the pretrained VLA was learned without
robot-action data.

\subsection{Flow Matching with Target Trajectories}
\label{sec:bg_flow_matching}

Flow matching learns a vector field that transports a simple source
distribution toward a target data distribution~\cite{flowmatching}.
For trajectory prediction, let $z_1$ collect the non-origin waypoints
of a target path and let $z_0$ be an independent standard Gaussian
sample of the same dimension. A common linear interpolation is
\begin{equation}
    z_\tau = (1-\tau)z_0 + \tau z_1,
    \qquad \tau\sim\mathcal{U}[0,1].
\end{equation}
The conditional vector field is trained with
\begin{equation}
    \mathcal{L}_{\mathrm{FM}}(\theta)
    = \mathbb{E}\!\left[
        \left\|v_\theta(z_\tau,\tau,I,g)-(z_1-z_0)\right\|_F^2
      \right],
\end{equation}
where the expectation is over observation-goal pairs, target
trajectories, source noise, and interpolation times. At inference,
integrating the learned field from $\tau=0$ to $\tau=1$ produces a
trajectory sample. This supervised construction requires target
trajectory samples. An energy gradient at a candidate path is a
different supervision signal; substituting it for $z_1-z_0$ does
not by itself retain the same distribution-matching objective.

\subsection{Geometric Energy and Local Refinement}
\label{sec:bg_energy}

A local Euclidean Signed Distance Field (ESDF) $\Phi$ assigns signed
distances to positions in the navigation plane. In the visibility-aware
representation used here, positive values indicate observed free space,
while occupied and unobserved regions receive non-positive values.
Interpolating this field along a
candidate path allows geometric preferences to be expressed through
a trajectory energy $E(x;g,\Phi)$. For language or image goals,
geometric goal terms use the corresponding training-time metric goal
annotation; this annotation is not an additional input to the learned
field unless waypoint conditioning is explicitly selected.

Where the energy is differentiable and its gradient is nonzero,
$-\nabla_x E(x;g,\Phi)$ gives an unconstrained local descent
direction. A projected gradient step takes the form
\begin{equation}
    x^{+}
    = \Pi_{\mathcal{K}}\!\left(
        x-\alpha\nabla_x E(x;g,\Phi)
      \right),
\end{equation}
where $\alpha>0$ is a step size and $\Pi_{\mathcal{K}}$ enforces
the chosen waypoint constraints, including the fixed origin.
Finite steps and constraint handling require care: an arbitrary
projected proposal need not decrease the energy. GEAR therefore
accounts for the acceptance of projected proposals when constructing
oracle labels, as detailed in Sec.~\ref{sec:oracle}. Waypoint
constraints alone do not guarantee obstacle avoidance.

\subsection{Dataset Aggregation for Sequential Prediction}
\label{sec:bg_aggregation}

When a predictor's output determines its next input, prediction
errors can move subsequent inputs away from the training
distribution. DAgger addresses this issue by collecting states
encountered by the current policy, obtaining expert labels at those
states, and retraining on their union with earlier
data~\cite{dagger}. In trajectory refinement, the analogous state
is a candidate path rather than a newly executed robot state.
GEAR applies this principle in trajectory space, using an
analytic energy oracle to label corrections. Its aggregation
mechanism targets the states encountered during refinement,
rather than only the initial trajectory distribution.


%%%%%%%%%%%%%%%%%%%%%%%%%%%%%%%%%%%%%%%%%%%%%%%%%%%%%%%%%%%%%%%%%%%%%%%%%%%%%%%%
\section{METHODOLOGY}
\label{sec:method}

\begin{figure*}[!t]
    \centering
    \includegraphics[width=1.0\textwidth]{fig/astar_pipeline.pdf}
    % Temporary caption; expand when the methodology text is finalized.
    \caption{GEAR training pipeline. Local energy supervision supplies correction targets, which are aggregated over candidate paths before continuation on additional training conditions. In full-dataset training, the current GEAR model generates candidate paths, the energy oracle supplies correction targets, and the same model learns from those targets. The updated model generates paths for the next cycle. The positive-part operator is $[z]_+=\max(0,z)$.}
    \label{fig:pipeline}
\end{figure*}

We introduce \textbf{GEAR}, an approach for training goal-conditioned navigation VLAs without target trajectory supervision. Rather than first solving a trajectory optimization problem and using the resulting trajectory as a training label, GEAR uses reconstructed scene geometry to define an analytic energy over candidate trajectories. The local descent direction of this energy provides privileged supervision for learning an autonomous trajectory-refinement field.
Figure~\ref{fig:pipeline} summarizes the training pipeline.

Given an RGB observation $I$, a goal specification $g$, and a candidate
waypoint trajectory
\[
    x = \{p_h\}_{h=0}^{H-1}, \qquad p_h \in \mathbb{R}^2,
\]
the learned field predicts a direction
\[
    F_\theta(I,g,x) \in \mathbb{R}^{H\times 2}
\]
that improves the current trajectory. During training, an ESDF $\Phi$ is
available as privileged geometric information. At deployment, the ESDF and
the analytic energy are discarded, and the learned field is iteratively
applied using only the visual observation, goal conditioning, and current
candidate trajectory.



\begin{figure}[!t]
    \centering
    \includegraphics[width=\columnwidth]{fig/model_simple.pdf}
    \caption{\textbf{GEAR correction model and refinement.}
      (a) RGB/goal context and the current path enter the Waypoint Action
      Expert, shown as one block with trainable adapters and correction head
      that predict $d^{(\ell)}$.
      Dashed branches regress energy-oracle corrections at aggregated paths,
      without target trajectories. Half-MSE excludes the fixed origin.
      The update caps correction magnitude, applies waypoint weights $A$, and
      enforces origin/segment bounds through $\Pi_{\mathcal K}$; the next path
      re-enters the same model.
      (b) Recorded initial, intermediate, and refined paths at a shared metric scale; dashed gray
      paths show the initial candidate. The ESDF is a visualization reference.
      Inference uses no ESDF or energy queries.}
    \label{fig:overview}
\end{figure}

\subsection{Action-Free Navigation Conditions}
\label{sec:conditions}

We construct training conditions from in-the-wild egocentric observations.
Each condition consists of an RGB image $I$, a navigation goal $g$, and a
local Euclidean Signed Distance Field (ESDF) $\Phi$ reconstructed from scene
geometry. Goals can be specified through language, an image region, or a
robot-frame waypoint.

Importantly, no reference navigation trajectory is associated with a training
condition. In particular, planned trajectories, robot demonstrations, and
trajectory actions are not used as supervision. The geometric representation
is used only to evaluate candidate trajectories and construct the local
energy-descent supervision described below.

For a waypoint $p$, we define robot-footprint clearance as
\[
    c(p) = \Phi(p) - r_{\mathrm{robot}},
\]
where $r_{\mathrm{robot}}$ is the robot radius. Occupied, unknown, invalid,
and out-of-map regions are treated conservatively when evaluating candidate
trajectories.

\subsection{Geometric Trajectory Energy}
\label{sec:energy}

We define a differentiable trajectory energy
\[
    E(x;g,\Phi)
    =
    \sum_{q\in\mathcal{Q}} \lambda_q E_q(x;g,\Phi),
\]
where the energy terms
\[
    \begin{gathered}
    \mathcal{Q}
    =
    \{
    \mathrm{collision},
    \mathrm{clearance},
    \mathrm{goal},\\
    \mathrm{progress},
    \mathrm{heading},
    \mathrm{smoothness}
    \}
    \end{gathered}
\]
capture complementary properties of a desirable navigation trajectory.

The collision term penalizes trajectories that violate robot-footprint
clearance or enter unsupported geometric regions. The clearance term can
encourage separation from nearby obstacles, while the goal and progress terms
encourage the trajectory to move toward the requested target. An early-heading
term discourages initially moving away from the goal, and a smoothness term
regularizes abrupt changes along the trajectory. Safety-related terms are
evaluated on dense samples along waypoint segments rather than only at the
predicted waypoints.

Unlike planner-based supervision, we do not minimize $E$ offline to obtain a
target path. Instead, the energy defines how the current candidate trajectory
should locally change.

\subsection{Projection-Aware Energy Oracle}
\label{sec:oracle}

For a candidate trajectory $x_k$ at aggregation depth $k$, we query an analytic
energy oracle by differentiating the trajectory energy with respect to the
trajectory itself:
\[
    g_k
    =
    \nabla_x E(x_k;g,\Phi).
\]

We convert the raw gradient into a normalized descent direction. Let
$\rho(g_k)$ denote the root-mean-square gradient magnitude over all non-origin
waypoints. The initial oracle direction is
\[
    \tilde d_k
    =
    -\frac{g_k}
    {\max\left(\rho(g_k),\epsilon_{\mathrm{grad}}\right)}.
\]
Directions below a minimum gradient magnitude are set to zero, and large
directions are capped to prevent unstable trajectory updates.

Candidate updates must additionally satisfy the trajectory parameterization
and kinematic constraints. Let $\Pi_{\mathcal K}$ denote the projection that
keeps the first waypoint fixed at the robot origin and converts the candidate
into a bounded, physically valid waypoint sequence. Starting from a nominal
step size $\eta$, we evaluate
\[
    x_k'
    =
    \Pi_{\mathcal K}
    \left(
        x_k + \alpha \tilde d_k
    \right),
\]
and geometrically backtrack $\alpha$ until a finite, non-energy-increasing
proposal is found:
\[
    E(x_k';g,\Phi)
    \leq
    E(x_k;g,\Phi) + \epsilon_E.
\]

If $\alpha_k$ is the accepted step size, the cached oracle label is the
correspondingly scaled pre-projection direction
\[
    d_k^\star
    =
    \frac{\alpha_k}{\eta}\tilde d_k.
\]
This construction makes the oracle supervision consistent with the projected
update that is subsequently applied to candidate trajectories.

\subsection{Autonomous VLA Trajectory Field}
\label{sec:field}



Figure~\ref{fig:overview} shows the correction model and iterative refinement procedure.

GEAR builds on a pretrained $\pi_{0.5}$ VLA backbone. Visual and goal
conditioning form the prefix context, while the current candidate waypoint
trajectory is provided to the action expert. The action expert is repurposed
to predict a trajectory-improvement field
\[
    F_\theta(I,g,x_k)
    \approx
    d_k^\star.
\]

In contrast to standard flow-matching training, we do not construct an
interpolation between noise and a ground-truth trajectory. The model time
token is fixed to a constant value throughout training and inference.
Consequently, $F_\theta$ represents an \emph{autonomous} vector field over
trajectory space: its prediction depends on the observation, goal, and
current trajectory, but not on aggregation depth $k$.

For an aggregated training tuple
$(I,g,x_k,d_k^\star)$, the field is trained using
\[
    \mathcal{L}_{\mathrm{field}}(\theta)
    =
    \frac{1}{2}
    \left\|
        F_\theta(I,g,x_k)_{1:H-1}
        -
        d_{k,1:H-1}^{\star}
    \right\|_2^2.
\]
The first waypoint is excluded because it is fixed at the robot origin.
The ESDF and energy computation do not appear inside this supervised SGD
objective; they are used only when querying the oracle.

\subsection{Energy-Based Dataset Aggregation}
\label{sec:aggregation}

Training the field only around a fixed initialization would introduce
distribution shift: after several learned refinement steps, the model may
encounter trajectories that are substantially different from the initial
training states. We therefore train GEAR using an iterative dataset
aggregation procedure.

For every training condition, we first sample a structured trajectory
$x_0$ from a goal-biased forward unicycle prior:
\[
    x_0
    \sim
    q_{\mathrm{prior}}(x\mid g).
\]
Unlike Gaussian noise, this initialization provides a coarse,
kinematically plausible trajectory that approximately points toward the goal.

At aggregation round $k$, we perform the following steps:
\begin{enumerate}
    \item Query the geometric oracle once at the current trajectory $x_k$ to
    obtain $d_k^\star$.
    \item Add the state-direction tuple
    $(I,g,x_k,d_k^\star)$ to the replay dataset.
    \item Train the VLA field on samples drawn from all previously aggregated
    rounds.
    \item Apply the newly trained field to the current frontier to obtain the
    next model-induced trajectory state.
\end{enumerate}

After $K$ rounds, the replay distribution is
\[
    \mathcal{D}_K
    =
    \bigcup_{k=0}^{K-1}
    \left\{
        (I_i,g_i,x_{i,k},d_{i,k}^\star)
    \right\}_{i=1}^{N}.
\]

During \emph{aggregation only}, the learned update is protected by the same
projection-aware energy backtracking used to validate an oracle step. Thus,
the next training frontier is obtained from the learned field while preventing
early inaccurate models from moving the replay distribution into
high-energy or invalid regions. The oracle is then queried again only after
this new frontier has been reached.

Aggregation continues until the trajectory updates converge or improvement
on a fixed, clip-disjoint model-only validation set plateaus. After collection
stops, the cached replay can optionally be used for additional SGD without
querying the oracle again.

\subsection{Model-Only Trajectory Refinement at Inference}
\label{sec:inference}

The analytic energy guard used during aggregation is strictly a
training-time mechanism. At inference, GEAR first constructs the same
goal-biased trajectory prior $x^{(0)}$ and repeatedly applies the learned
autonomous field:
\[
    \begin{gathered}
    x^{(\ell+1)}
    =
    \Pi_{\mathcal K}
    \left(
        x^{(\ell)}
        +
        \eta
        F_\theta(I,g,x^{(\ell)})
    \right),\\
    \ell=0,\ldots,L-1.
    \end{gathered}
\]

We use a fixed refinement budget of $L=12$ steps in our experiments.
No ESDF query, trajectory-energy evaluation, negative-energy gradient, or
energy-based acceptance test is performed during this rollout. The resulting
policy therefore requires only the visual observation, goal specification,
and lightweight kinematic projection at deployment.


%%%%%%%%%%%%%%%%%%%%%%%%%%%%%%%%%%%%%%%%%%%%%%%%%%%%%%%%%%%%%%%%%%%%%%%%%%%%%%%%
\section{ANALYSIS AND RESULTS}

\label{sec:results}

\subsection{Experimental Setup}
\label{sec:experimental_setup}

\textbf{Baselines.}
We adopt the same three navigation baselines as VEGA~\cite{vega}:
OmniVLA~\cite{omnivla}, an omni-modal navigation VLA supporting language,
image, and spatial goals; TIC-VLA~\cite{ticvla}, a think-in-control VLA for
navigation in dynamic environments; and NaVILA~\cite{navila}, a
vision-language-action model for legged robot navigation.

GEAR-Nav's task-level evaluation will use the \emph{model-only} rollout
defined in Sec.~\ref{sec:inference}, without the energy oracle or the
training-time acceptance guard.

\textbf{Evaluation protocol.}
Our comparison covers Isaac-Sim navigation and real-world robot navigation.
The baseline values in Tables~\ref{tab:isaac_sim} and~\ref{tab:real_world}
are reproduced from VEGA~\cite{vega}; they are not new measurements from
our evaluation. The GEAR-Nav simulation values in Table~\ref{tab:isaac_sim}
were supplied separately; their run provenance and comparability with the
reference protocol remain to be verified. Real-world GEAR-Nav results are
pending. In the reference real-world evaluation,
policy outputs use a shared waypoint-tracking stack without additional path
planning or obstacle avoidance. NaVILA's high-level commands are converted
to local waypoints; its native LiDAR-conditioned obstacle-avoiding locomotion
policy is not used~\cite{vega}.

\textbf{Metrics.}
For the simulation and real-world comparisons, we retain the metrics and
units reported in VEGA~\cite{vega}:
\begin{itemize}
    \item \textbf{Success rate ($S$)}: fraction of successful goal-reaching
    trials. In the real-world evaluation, success requires reaching the goal
    without collision or safety intervention.
    \item \textbf{Collision rate ($C$)}: fraction of trials with a collision;
    real-world safety interventions also count as collisions.
    \item \textbf{Obstacle clearance ($Cl$)}: average minimum distance to
    obstacles over non-colliding trials, in meters. N/A denotes cases in
    which all trials collide.
\end{itemize}
Success and collision rates are reported on a $[0,1]$ scale.

For the refinement analysis in Sec.~\ref{sec:energy_analysis}, we additionally
consider goal progress, required-progress ratio, dense collision rate, and
safe success under the robot-footprint margin. These geometric diagnostics
are distinct from the trial-level metrics in the comparison tables.

\subsection{Simulation Evaluation}
\label{sec:simulation}

We adopt three Isaac-Sim scenarios containing static obstacles and moving
pedestrians. The reference evaluation varies pedestrian motion onset using
identical timing conditions across methods, with 50 trials per method and
scenario~\cite{vega}. Table~\ref{tab:isaac_sim} records the reported baseline
results and the supplied GEAR-Nav values.

\begin{table}[!tb]
\centering
\caption{Isaac-Sim navigation. VEGA baselines use 50 trials per scenario.}
\label{tab:isaac_sim}
\footnotesize
\setlength{\tabcolsep}{2pt}
\renewcommand{\arraystretch}{1.08}
\begin{tabular*}{\columnwidth}{@{\extracolsep{\fill}}lccc@{\hspace{7pt}}ccc@{\hspace{7pt}}ccc@{}}
\toprule
& \multicolumn{3}{c}{\textbf{Scenario 1}}
& \multicolumn{3}{c}{\textbf{Scenario 2}}
& \multicolumn{3}{c}{\textbf{Scenario 3}} \\
\cmidrule(lr){2-4}\cmidrule(lr){5-7}\cmidrule(lr){8-10}
\textbf{Method}
& $S\uparrow$ & $C\downarrow$ & $Cl\uparrow$
& $S\uparrow$ & $C\downarrow$ & $Cl\uparrow$
& $S\uparrow$ & $C\downarrow$ & $Cl\uparrow$ \\
\midrule
OmniVLA
& 0.42 & 0.58 & 0.04
& 0.06 & 0.96 & 0.00
& 0.04 & 0.16 & 0.13 \\
TIC-VLA
& 0.18 & 0.50 & 0.45
& 0.90 & 0.04 & 0.48
& 0.74 & 0.06 & 0.61 \\
NaVILA
& 0.80 & 0.00 & 0.84
& 0.90 & 0.02 & 0.53
& 0.00 & 0.00 & 1.52 \\
\midrule
\textbf{GEAR-Nav}
& 0.92 & 0.06 & 0.78
& 0.84 & 0.10 & 0.46
& 0.86 & 0.12 & 0.71 \\
\bottomrule
\end{tabular*}
\par\vspace{2pt}
\parbox{\columnwidth}{\scriptsize\raggedright
$S$: success rate; $C$: collision rate; $Cl$: clearance (m).\\
Baselines: VEGA~\cite{vega}. GEAR-Nav values supplied separately; run provenance and protocol comparability require verification.}
\end{table}

\subsection{Real-World Robot Evaluation}
\label{sec:real_world}

Our real-world comparison covers three cluttered scenarios containing static
and dynamic obstacles. Table~\ref{tab:real_world} reproduces the baseline
results for Scenarios~1--3 of VEGA's real-world evaluation, with 20 trials
per method and scenario~\cite{vega}. GEAR-Nav's quantitative results remain
pending.
Figure~\ref{fig:real_world_qualitative} illustrates schematic routes in these
scenarios. GEAR-Nav routes summarize our recordings; the baseline overlays
are adapted from the qualitative VEGA comparison~\cite{vega}.

\begin{table}[!tb]
\centering
\caption{Real-world navigation (20 trials per method and scenario).}
\label{tab:real_world}
\footnotesize
\setlength{\tabcolsep}{2pt}
\renewcommand{\arraystretch}{1.08}
\begin{tabular*}{\columnwidth}{@{\extracolsep{\fill}}lccc@{\hspace{7pt}}ccc@{\hspace{7pt}}ccc@{}}
\toprule
& \multicolumn{3}{c}{\textbf{Scenario 1}}
& \multicolumn{3}{c}{\textbf{Scenario 2}}
& \multicolumn{3}{c}{\textbf{Scenario 3}} \\
\cmidrule(lr){2-4}\cmidrule(lr){5-7}\cmidrule(lr){8-10}
\textbf{Method}
& $S\uparrow$ & $C\downarrow$ & $Cl\uparrow$
& $S\uparrow$ & $C\downarrow$ & $Cl\uparrow$
& $S\uparrow$ & $C\downarrow$ & $Cl\uparrow$ \\
\midrule
OmniVLA
& 0.15 & 0.85 & 0.35
& 0.00 & 0.65 & 0.32
& 0.00 & 0.50 & 0.44 \\
TIC-VLA
& 0.00 & 1.00 & N/A
& 0.10 & 0.85 & 0.37
& 0.00 & 1.00 & N/A \\
NaVILA
& 0.00 & 1.00 & N/A
& 0.00 & 0.80 & 0.35
& 0.00 & 0.45 & 0.31 \\
\midrule
\textbf{GEAR-Nav}
& -- & -- & --
& -- & -- & --
& -- & -- & -- \\
\bottomrule
\end{tabular*}
\par\vspace{2pt}
\parbox{\columnwidth}{\scriptsize\raggedright
$S$: success rate; $C$: collision rate; $Cl$: clearance (m).\\
Baselines: VEGA~\cite{vega}. ``--'': pending; N/A: all trials collide.}
\end{table}

\begin{figure}[!tb]
    \centering
    \includegraphics[width=\columnwidth]{fig/astar_realworld.pdf}
    \refstepcounter{figure}
    \label{fig:real_world_qualitative}
\end{figure}

The real-world experiments will test whether the geometric behavior distilled
through energy-based aggregation transfers to scenes where no ESDF or
trajectory-energy computation is available to the deployed policy. In
particular, we will examine whether GEAR can remain goal-directed while
modifying its trajectory around intervening static obstacles and newly
observed dynamic obstacles.

\subsection{Analysis of Learned Energy Refinement}
\label{sec:energy_analysis}

We next isolate whether the learned field actually improves the structured
trajectory prior and how closely it approaches privileged geometric
optimization. We compare four rollout modes:
\begin{itemize}
    \item \textbf{Prior}: the goal-biased unicycle initialization without
    learned refinement.
    \item \textbf{Model-only}: the deployable GEAR rollout using only the
    learned field and kinematic projection.
    \item \textbf{Guarded}: the learned field combined with privileged
    analytic-energy backtracking.
    \item \textbf{Oracle}: direct refinement using the analytic
    negative-energy-gradient oracle.
\end{itemize}

Guarded and Oracle use privileged scene geometry and are included only as
diagnostic references. These comparisons will evaluate the final trajectory
after the fixed refinement budget.

The comparison between Prior and Model-only measures the improvement provided
by the learned trajectory field itself. The difference between Model-only and
Oracle characterizes the remaining gap between the visual policy and the
privileged analytic optimizer, while the Guarded result indicates how much of
this gap is caused by occasional inaccurate learned update directions or step
magnitudes.

We additionally analyze performance as a function of refinement depth. A
successful autonomous field should progressively transform the structured
prior into trajectories with greater goal progress and lower collision rate,
without requiring access to the geometric energy during inference.

% TODO: Insert plot:
% x-axis: model-only refinement depth 0,...,12
% y-axis: safe success / collision rate / progress ratio
% optionally overlay oracle as a privileged reference.

\balance

%%%%%%%%%%%%%%%%%%%%%%%%%%%%%%%%%%%%%%%%%%%%%%%%%%%%%%%%%%%%%%%%%%%%%%%%%%%%%%%%



%%%%%%%%%%%%%%%%%%%%%%%%%%%%%%%%%%%%%%%%%%%%%%%%%%%%%%%%%%%%%%%%%%%%%%%%%%%%%%%%



%%%%%%%%%%%%%%%%%%%%%%%%%%%%%%%%%%%%%%%%%%%%%%%%%%%%%%%%%%%%%%%%%%%%%%%%%%%%%%%%
\section*{APPENDIX}


\section*{ACKNOWLEDGMENT}



%%%%%%%%%%%%%%%%%%%%%%%%%%%%%%%%%%%%%%%%%%%%%%%%%%%%%%%%%%%%%%%%%%%%%%%%%%%%%%%%


\begin{thebibliography}{99}
    \bibitem{vega}
    G. Seneviratne, Y. Abeysinghe, J. An, V. Shende, R. Kumar, and
    D. Manocha, ``VEGA: Learning Navigation VLAs from In-the-Wild
    Egocentric Video with Geometric Trajectory Supervision,'' manuscript.

    \bibitem{omnivla}
    N. Hirose, C. Glossop, D. Shah, and S. Levine,
    ``OmniVLA: An Omni-Modal Vision-Language-Action Model for Robot
    Navigation,'' arXiv:2509.19480, 2025.

    \bibitem{ticvla}
    Z. Huang, Y. Zhang, J. Liu, R. Song, C. Tang, and J. Ma,
    ``TIC-VLA: A Think-in-Control Vision-Language-Action Model for Robot
    Navigation in Dynamic Environments,'' arXiv:2602.02459, 2026.

    \bibitem{navila}
    A.-C. Cheng \emph{et al.}, ``NaVILA: Legged Robot
    Vision-Language-Action Model for Navigation,'' arXiv:2412.04453, 2025.

\end{thebibliography}




\ifincludesupplementary
  \clearpage
  \nobalance
  \onecolumn
  \input{supplementary}
\fi

\end{document}
